\origpage{342}{N. Robertson, P.D. Seymour / Journal of Combinatorial Theory, Series B 90 (2004) 325--385}
\noindent that $rep(G'\setminus V(Q))\ge 1$, since $k$ satisfies 6.5 unless $\Sigma$ is a sphere. Let $F$ be an $O$-arc in $\Sigma$ with $U(Q)\subseteq F$ and $F\setminus U(Q)\subseteq r$. Since $F\cap U(G'\setminus V(Q)) = \emptyset$ and $rep(G'\setminus V(Q))\ge 1$, it follows that there is a disc $\Delta\subseteq \Sigma$ bounded by $F$.

\hspace*{12pt}(1) \textit{If $G'\cap \Delta\subseteq \sigma'(v)$ then there is a path $P$ of $\sigma'(v)$ with ends $a, b$ and with $U(P)\subseteq \bar{r}$.}

\textbf{Subproof.} $r\cap F$ is connected, and so there is a path of $G'\cap (\bar{r}\cap \Delta)$ between $a$ and $b$; and so if $G'\cap \Delta\subseteq \sigma'(v)$ we may take $P$ to be this path. This proves (1).

\hspace*{12pt}(2) \textit{If $U(\sigma'(u))\subseteq \Delta$ for every $u\in V(G)\setminus \{v\}$ then there is a path $P$ of $\sigma'(v)$ with ends $a, b$ and with $U(P)\subseteq \bar{r}$.}

\textbf{Subproof.} Let $G''$ be $G'\cap \Delta$, regarded as a drawing in $\Sigma$. If $\Sigma$ is not a sphere then there is a non-null-homotopic $O$-arc disjoint from $\Delta$, and hence $rep(G'') = 0$, which is impossible under the hypothesis of (2), because there is an enlargement of $H$ in $G''$, and $rep(H)\ge k$, and $k$ satisfies 6.5. Thus $\Sigma$ is a sphere. Let $\Delta'\subseteq \Sigma$ be the disc bounded by $F$ with $\Delta\ne \Delta'$. Then $G'\cap \Delta'\subseteq \sigma'(v)$, since $U(\sigma'(u))\subseteq \Delta$ for every $u\in V(G)\setminus \{v\}$ and the second statement of the theorem holds. But then the claim follows by (1) applied to $\Delta'$. This proves (2).

From (1) and (2) we may suppose, for a contradiction, that there exist $u_1, u_2\in V(G)\setminus \{v\}$ with $U(\sigma'(u_1))\cap \Delta\ne \emptyset$ and $U(\sigma'(u_2))\nsubseteq \Delta$. Since $U(\sigma'(u_i))\cap F = \emptyset$ and $\sigma'(u_i)$ is connected ($i = 1, 2$), it follows that $U(\sigma'(u_1))\subseteq \Delta\setminus F$ and $U(\sigma'(u_2))\cap \Delta = \emptyset$. In particular, $u_1\ne u_2$. Now every path of $G'$ from $V(\sigma'(u_1))$ to $V(\sigma'(u_2))$ passes through $V(\sigma'(v))$, and so every path of $G$ from $u_1$ to $u_2$ passes through $v$. But $G$ is a block, a contradiction, as required. $\square$

\fix{We prove the fact used in the proof of 6.6, that if $G$ is a drawing in a surface $\Sigma$ with $bd(\Sigma) = \emptyset$ which is not a sphere, and $v \in V(G)$, then $rep(G) \le rep(G \setminus v) + 1$. We need the following.

\hspace*{12pt}$(*)$ \textit{Let $D \subseteq \Sigma$ be a closed disc, and let $F$ be a non-null-homotopic $O$-arc with $F \cap D \ne \emptyset$. For all $a, b \in bd(D)$ let $P(a, b) \subseteq D$ be an arc with ends $a, b$ (with $P(a, a) = \{a\}$). Then there is a component $A$ of $F \setminus D$, with ends $a, b$ say, such that $A \cup P(a, b)$ is a non-null-homotopic $O$-arc.}

\textbf{Subproof.} Let $D'$ be a closed disc with $D$ in its interior. Since $F$ is compact, only finitely many components of $F \setminus D$ meet $\Sigma \setminus D'$, say $A_1, \dots, A_m$, where $A_i$ has ends $a_i, b_i$; and $F$ is the union of $A_1, \dots, A_m$ and finitely many arcs of $F$ in $D'$, which join them up in a cyclic order. Suppose that $A_i \cup P(a_i, b_i)$ is null-homotopic for $1 \le i \le m$. Then each $A_i$ is homotopic, with its ends fixed, to $P(a_i, b_i)$, which lies in $D'$. Hence $F$ is homotopic to a closed curve in $D'$, and so $F$ is null-homotopic, since $D'$ is a disc, a contradiction. So some $A_i \cup P(a_i, b_i)$ is non-null-homotopic. This proves $(*)$.

Now let $H = G \setminus v$, and let $F_0$ be a non-null-homotopic $H$-normal $O$-arc with $|F_0 \cap V(H)| = rep(H)$. Let $e_1, \dots, e_t$ be the edges of $G$ incident with $v$, and for $0 \le i \le t$ let $H_i$ be the drawing obtained from $H$ by adding $v$ and $e_1, \dots, e_i$. We show by induction on $i$ that there is a non-null-homotopic $O$-arc $F_i$ with $F_i \cap U(H_i) \subseteq (F_0 \cap V(H)) \cup \{v\}$. For $i = 0$ this holds since $F_0$ is $H$-normal. Let $1 \le i \le t$, let $u$ be the end of $e_i$ different from $v$, and let $Z = U(H_{i-1})$. Since $Z$ is closed and $\bar{e}_i \cap Z = \{u, v\}$, there is a closed disc $D$, a thin neighbourhood of $\bar{e}_i$, with $\bar{e}_i \subseteq D$, $u, v \in bd(D)$, $bd(D) \cap \bar{e}_i = \{u, v\}$ and $D \cap Z = \{u, v\}$. If $F_{i-1} \cap D = \emptyset$, let $F_i = F_{i-1}$. Otherwise, for $a, b \in bd(D)$ let $P(a, b)$ be an arc of $bd(D)$ with ends $a, b$ that does not contain $u$ unless $u \in \{a, b\}$, which exists since $bd(D)$ is an $O$-arc; and let $F_i = A \cup P(a, b)$, where $A, a, b$ are as in $(*)$ applied to $F_{i-1}$. Then $F_i$ is non-null-homotopic. Since $A \subseteq F_{i-1} \setminus D$ and $\bar{e}_i \subseteq D$, we have $A \cap U(H_i) \subseteq F_{i-1} \cap Z$. Moreover $P(a, b) \cap U(H_i) \subseteq bd(D) \cap (Z \cup \bar{e}_i) = \{u, v\}$, and if $u \in P(a, b)$ then $u \in \{a, b\} \subseteq F_{i-1}$. Hence $F_i \cap U(H_i) \subseteq (F_{i-1} \cap Z) \cup \{v\} \subseteq (F_0 \cap V(H)) \cup \{v\}$, which completes the induction. Since $H_t = G$, the $O$-arc $F_t$ is non-null-homotopic, and $F_t \cap U(G) \subseteq (F_0 \cap V(H)) \cup \{v\} \subseteq V(G)$. Thus $F_t$ is $G$-normal, and $rep(G) \le |F_t \cap V(G)| \le rep(H) + 1$, as required.}\fixnote{fix:6.6-rep-proof-added}{The paper gives no reason why $rep(H) \ge k$, that is, why deleting a vertex lowers $rep$ by at most one. The argument above is the one used in the formalization.}

\textbf{6.7.} \textit{If $\Gamma$ is an internally 3-connected painting in a surface $\Sigma$, then $sk(\Gamma)$ is a block and $rep(sk(\Gamma)) = rep(\Gamma)$.}

\textbf{Proof.} We may assume that $|V(\Gamma)|\ge 3$. Now since for every region $r$ of $\Gamma$ in $\Sigma$, $r$ is an open disc and $r\cup \{v\}$ is simply-connected for every vertex $v$ of $\Gamma$ with $v\in \bar{r}$, it follows that the closed curve tracing the perimeter of $r$ has no repeated vertices, and since $|V(\Gamma)|\ge 3$ this curve is an $O$-arc. Thus for each region $r$ of $\Gamma$ in $\Sigma$ there is a circuit $C_r$ of $sk(\Gamma)$ with $U(C_r) = bd(\bar{r})$. Hence the same is true for each region $r$ of $sk(\Gamma)$ in $\hat{\Sigma}$. Now if $v\in V(sk(\Gamma))$ and $sk(\Gamma)\setminus v$ is disconnected, then there are two vertices $u_1, u_2\ne v$ belonging to the same circuit $C_r$, and belonging to different components of $sk(\Gamma)\setminus v$, which is impossible. Thus $sk(\Gamma)$ is a block. The second claim is clear. $\square$

\textbf{Proof of 6.2.} Choose $n$ such that 6.6 holds with $\Sigma$ and $k$ replaced by $\hat{\Sigma}$ and $n - 1$. Let $\Phi, \Gamma_1, \Gamma_2, \sigma$ be as in 6.2. For each $e\in E(sk(\Gamma_1))$ define $\tau(e)$ as follows. If $e\in E(\Gamma_1)$ let $\tau(e) = \sigma(e)$. If $e$ is a component of $bd(\bar{c}_1)\setminus \tilde{c}_1$ for some $c_1\in E(\Gamma_1)$ with $|\tilde{c}_1| = 3$, let $c_2 = \sigma(c_1)$, let the ends of $e$ be $u, v$ and let $\tau(e)$ be the edge of $sk(\Gamma_2)$ with $\tau(e)\subseteq bd(\bar{c}_2)$ and with ends in $\sigma(u), \sigma(v)$. For each $v\in V(sk(\Gamma_1))$ let $\tau(v) = \sigma(v)$. Then $\tau$ is an enlargement of $sk(\Gamma_1)$ in $sk(\Gamma_2)$.
